\section{变幅内摆线}
\markboth{35. 变幅内摆线}{35. 变幅内摆线}

设 $r<R$，当半径为 $r$ 的动圆 $O_1$ 在半径为 $R$ 的定圆 $O$ 内无滑动地滚动时，在圆 $O_1$ 平面内与圆 $O_1$ 固定连结但不在 $O_1$ 圆周上的一点 $M$ 的轨迹叫做变幅内摆线。

变幅内摆线可分成两类：当 $M$ 在圆 $O_1$ 外部时，曲线称为长幅内摆线；而当 $M$ 在圆 $O_1$ 内部时，称为短幅内摆线。

如果 $M$ 在圆 $O_1$ 的圆周上，那么 $M$ 的轨迹成为内摆线，这在前面第 24 节中已经讨论过。内摆线和变幅内摆线统称为内摆线族曲线。

例如，在第 22 节中讨论卡丹的转盘问题时，我们实际上已经证明了，当动圆半径等于定圆半径的一半时，长幅内摆线和短幅内摆线都是椭圆。

\noindent
\begin{minipage}[t]{0.44\textwidth}
  \vspace{0pt}
  \centering
  \begin{tikzpicture}[scale=0.84, >=Stealth, line cap=round, line join=round]
    % R = 1.52, r = 0.66, h = 0.40, angle(AON) = 52 deg
    \def\R{1.52}
    \def\r{0.66}
    \def\h{0.40}
    \def\angAON{52}
    \pgfmathsetmacro{\tdeg}{\R/\r*\angAON}
    \pgfmathsetmacro{\alp}{\angAON - \tdeg}
    \coordinate (O) at (0,0);
    \coordinate (N) at ({\R*cos(\angAON)}, {\R*sin(\angAON)});
    \coordinate (O1) at ({(\R-\r)*cos(\angAON)}, {(\R-\r)*sin(\angAON)});
    \coordinate (M) at ({(\R-\r)*cos(\angAON) + \h*cos(\alp)}, {(\R-\r)*sin(\angAON) + \h*sin(\alp)});

    \draw[->] (-1.70, 0) -- (1.88, 0) node[below, font=\footnotesize] {$x$};
    \draw[->] (0, -1.70) -- (0, 1.88) node[left, font=\footnotesize] {$y$};

    \draw[thick] (O) circle (\R);
    \draw (O1) circle (\r);

    \draw (O) -- (N);
    \draw (O1) -- (M);
    \draw[thin] (O1) -- ++(0.72, 0);

    % 角 t 与 角 alpha
    \draw ($(O1) + (\angAON:0.24)$) arc (\angAON:\alp:0.24);
    \node[font=\scriptsize] at ($(O1) + (26:0.38)$) {$t$};
    \draw ($(O1) + (0:0.32)$) arc (0:\alp:0.32);
    \node[font=\scriptsize] at ($(O1) + ({0.5*\alp}:0.48)$) {$\alpha$};

    \fill (O) circle (1.2pt) node[below left, font=\footnotesize] {$O$};
    \fill (N) circle (1.2pt) node[above right=-1pt, font=\footnotesize] {$N$};
    \fill (O1) circle (1.2pt) node[above left=-1pt, font=\footnotesize] {$O_1$};
    \fill (M) circle (1.3pt) node[right=2pt, font=\footnotesize] {$M$};
  \end{tikzpicture}
  \captionof*{figure}{图\ 35.1}
\end{minipage}
\hfill
\begin{minipage}[t]{0.53\textwidth}
  \vspace{0pt}
  \hspace{2em}下面推导变幅内摆线的方程。如图 35.1，取定圆的圆心 $O$ 为坐标原点，并且选取 $x$ 轴的方向，使动圆的圆心落在 $x$ 轴正半轴上时，$M$ 的位置离 $O$ 点最远。设动圆 $O_1$ 滚动到图 35.1 中的任意位置，这时 $M$ 的坐标是 $(x, y)$，并且两圆的切点是 $N$。仿 24 节内摆线的求法，设
  \[
    \angle MO_1N = t,
  \]
\end{minipage}

\noindent
则 $\overrightarrow{O_1M}$ 的幅角 $\alpha=\frac{r}{R}t-t$，记 $O_1M=h$，由 $\overrightarrow{OM}=\overrightarrow{OO_1}+\overrightarrow{O_1M}$ 得
\begin{align*}
  (x, y) &= \Bigl((R-r)\cos\frac{r}{R}t,\, (R-r)\sin\frac{r}{R}t\Bigr) + (h\cos\alpha,\, h\sin\alpha) \\
         &= \biggl((R-r)\cos\frac{r}{R}t + h\cos\Bigl(t - \frac{r}{R}t\Bigr), \\
         &\qquad\ (R-r)\sin\frac{r}{R}t - h\sin\Bigl(t - \frac{r}{R}t\Bigr)\biggr).
\end{align*}
令 $\frac{r}{R}=m$，就得到曲线的参数方程为
\[
  \begin{aligned}
    x &= (R-mR)\cos mt + h\cos(t-mt), \\
    y &= (R-mR)\sin mt - h\sin(t-mt).
  \end{aligned}
  \quad (0<m<1)
\]
这里，由于 $0<r<R$，所以 $0<m<1$。当 $h>mR$ 时得到长幅内摆线，$h<mR$ 时得到短幅内摆线。若 $h=mR$，则得到内摆线。所以上述方程是内摆线族曲线的统一方程。

在图 35.2～35.4 中分别画出了 $m=\frac{1}{4}$ 时的长幅内摆线、内摆线（星形线）和短幅内摆线。

\begin{figure}[htbp]
  \centering
  \begin{minipage}[b]{0.32\textwidth}
    \centering
    \begin{tikzpicture}[scale=1.08, line cap=round, line join=round]
      % m = 1/4, R = 1, r = 0.25, h = 0.45 (长幅内摆线)
      \def\R{1.0}
      \def\r{0.25}
      \def\h{0.45}
      \def\ang{38}
      \pgfmathsetmacro{\tdeg}{\R/\r*\ang}
      \pgfmathsetmacro{\alp}{\ang - \tdeg}
      \draw (0,0) circle (\R);
      \draw ({(\R-\r)*cos(\ang)}, {(\R-\r)*sin(\ang)}) circle (\r);
      \draw ({(\R-\r)*cos(\ang)}, {(\R-\r)*sin(\ang)}) -- ({(\R-\r)*cos(\ang)+\h*cos(\alp)}, {(\R-\r)*sin(\ang)+\h*sin(\alp)});
      \fill ({(\R-\r)*cos(\ang)}, {(\R-\r)*sin(\ang)}) circle (0.9pt);
      \fill ({(\R-\r)*cos(\ang)+\h*cos(\alp)}, {(\R-\r)*sin(\ang)+\h*sin(\alp)}) circle (1.1pt);
      \draw[very thick, curveblue, domain=0:1440, samples=260, variable=\s]
        plot ({(\R-\r)*cos(0.25*\s) + \h*cos(0.75*\s)},
              {(\R-\r)*sin(0.25*\s) - \h*sin(0.75*\s)});
    \end{tikzpicture}
    \caption*{图\ 35.2}
  \end{minipage}
  \hfill
  \begin{minipage}[b]{0.32\textwidth}
    \centering
    \begin{tikzpicture}[scale=1.08, line cap=round, line join=round]
      % m = 1/4, R = 1, r = 0.25, h = 0.25 (内摆线/星形线)
      \def\R{1.0}
      \def\r{0.25}
      \def\h{0.25}
      \def\ang{38}
      \pgfmathsetmacro{\tdeg}{\R/\r*\ang}
      \pgfmathsetmacro{\alp}{\ang - \tdeg}
      \draw (0,0) circle (\R);
      \draw ({(\R-\r)*cos(\ang)}, {(\R-\r)*sin(\ang)}) circle (\r);
      \draw ({(\R-\r)*cos(\ang)}, {(\R-\r)*sin(\ang)}) -- ({(\R-\r)*cos(\ang)+\h*cos(\alp)}, {(\R-\r)*sin(\ang)+\h*sin(\alp)});
      \fill ({(\R-\r)*cos(\ang)}, {(\R-\r)*sin(\ang)}) circle (0.9pt);
      \fill ({(\R-\r)*cos(\ang)+\h*cos(\alp)}, {(\R-\r)*sin(\ang)+\h*sin(\alp)}) circle (1.1pt);
      \draw[very thick, curveblue, domain=0:1440, samples=260, variable=\s]
        plot ({(\R-\r)*cos(0.25*\s) + \h*cos(0.75*\s)},
              {(\R-\r)*sin(0.25*\s) - \h*sin(0.75*\s)});
    \end{tikzpicture}
    \caption*{图\ 35.3}
  \end{minipage}
  \hfill
  \begin{minipage}[b]{0.32\textwidth}
    \centering
    \begin{tikzpicture}[scale=1.08, line cap=round, line join=round]
      % m = 1/4, R = 1, r = 0.25, h = 0.13 (短幅内摆线)
      \def\R{1.0}
      \def\r{0.25}
      \def\h{0.13}
      \def\ang{38}
      \pgfmathsetmacro{\tdeg}{\R/\r*\ang}
      \pgfmathsetmacro{\alp}{\ang - \tdeg}
      \draw (0,0) circle (\R);
      \draw ({(\R-\r)*cos(\ang)}, {(\R-\r)*sin(\ang)}) circle (\r);
      \draw ({(\R-\r)*cos(\ang)}, {(\R-\r)*sin(\ang)}) -- ({(\R-\r)*cos(\ang)+\h*cos(\alp)}, {(\R-\r)*sin(\ang)+\h*sin(\alp)});
      \fill ({(\R-\r)*cos(\ang)}, {(\R-\r)*sin(\ang)}) circle (0.9pt);
      \fill ({(\R-\r)*cos(\ang)+\h*cos(\alp)}, {(\R-\r)*sin(\ang)+\h*sin(\alp)}) circle (1.1pt);
      \draw[very thick, curveblue, domain=0:1440, samples=260, variable=\s]
        plot ({(\R-\r)*cos(0.25*\s) + \h*cos(0.75*\s)},
              {(\R-\r)*sin(0.25*\s) - \h*sin(0.75*\s)});
    \end{tikzpicture}
    \caption*{图\ 35.4}
  \end{minipage}
\end{figure}
